\documentclass[../readings.tex]{subfiles}

\begin{document}

\subsection{Fourier Analysis as Change of Basis}
\label{sub:fourier}

Diagonalizing the Laplacian via sinusoidal eigenfunctions. Turn differential equations into algebra.

\subsubsection{Discrete Fourier Transform (DFT)}

\addDef{DFT Matrix}{%
Unitary matrix $\mathbf{F} \in \fC^{n \times n}$ with entries:
\[
    F_{jk} = \frac{1}{\sqrt{n}} \omega_n^{jk}, \quad \omega_n = e^{-2\pi i/n}.
\]
\begin{itemize}
    \item \textbf{Unitary Property:}\enspace $\mathbf{F}^* \mathbf{F} = \bI$. Orthonormal columns.
    \item \textbf{Inverse DFT:}\enspace $\mathbf{F}^{-1} = \mathbf{F}^*$.
\end{itemize}
}

\addThm{Parseval's Theorem}{%
The DFT is an isometry: $\|\mathbf{F}\bx\|_2 = \|\bx\|_2$. Total energy is preserved between time and frequency domains.
}

\subsubsection{Convolution and Circulant Matrices}

\addDef{Circulant Matrix}{%
Matrix $\mathbf{C}$ where each row is a cyclic shift of the first column $\bc$.
\begin{itemize}
    \item \textbf{Diagonalization:}\enspace Every circulant matrix is diagonalized by the DFT: $\mathbf{C} = \mathbf{F}^* \text{diag}(\mathbf{F}\bc) \mathbf{F}$.
    \item \textbf{Eigenvalues:}\enspace The eigenvalues of $\mathbf{C}$ are the DFT coefficients of its first column.
\end{itemize}
}

\addThm{Convolution Theorem}{%
Circular convolution $\bz = \bc * \bx$ is equivalent to pointwise multiplication in frequency:
\[
    \mathbf{F}(\bc * \bx) = (\mathbf{F}\bc) \odot (\mathbf{F}\bx).
\]
\textbf{Complexity:}\enspace Reduced from $O(n^2)$ (direct) to $O(n \log n)$ (FFT).
}

\subsubsection{The Fast Fourier Transform (FFT)}

\addThm{Cooley-Tukey Algorithm}{%
Computes the DFT in \textbf{$O(n \log n)$} flops for $n = 2^p$.
\begin{itemize}
    \item \textbf{Mechanism:}\enspace Recursively splits $n$-point DFT into two $n/2$-point DFTs (even/odd indices).
    \item \textbf{NumPy:}\enspace Use \texttt{np.fft.fft} and \texttt{np.fft.ifft}.
\end{itemize}
}

\addNote{%
\textbf{Spectral Differentiation:}\enspace For smooth periodic signals, differentiate by multiplying frequency components $\hat{x}_k$ by $ik$. Provides exponential accuracy compared to polynomial accuracy of finite differences.
}

\subsubsection{Exercises}

\addExcr{%
\begin{enumerate}
    \item Write out the $4 \times 4$ DFT matrix explicitly. Verify it is unitary.
    \item Solve the periodic Poisson equation $-u'' = f$ via FFT. (\textit{Hint: Laplacian is circulant}).
    \item Compare the speed of \texttt{np.fft.fft(x)} vs.\ \texttt{F @ x} for $n=4096$.
    \item Implement spectral differentiation for $f(x) = \sin(x)$ and check error vs.\ $n$.
\end{enumerate}
}

\end{document}
